\section*{Part VIII: The Continuing Case: Bayes Institute's Economic Architecture}

\subsection*{76. The Institution Is an Ambition; the First Delivery System Is an Experiment}

Bayes Institute's enduring mission is to create the conditions in which human possibility becomes human capability. Its long-range ambition includes a serious STEM institution, research, industry engagement, and potentially physical and international educational infrastructure. None of these aspirations establishes that a particular tuition, 6-month program, partnership package, or enterprise contract will work now.

At its early stage, Bayes's most immediately accessible assets are founder expertise in technical subjects, experience teaching and designing programs, and the ability to create rigorous instructional work. Its constraints include limited founder hours alongside a full-time profession, scarce discretionary investment capital, the absence of an established institutional credential, and the need to earn trust without pretending that research infrastructure, faculty depth, placement outcomes, or alumni prestige already exist.

Consequently the delivery strategy should be \textbf{asset-light, founder-informed, and deliberately instrumented}. That does not mean every product must be cheap or every operation must be improvised. It means the institute should make each fixed commitment only when that commitment has an economically and educationally defensible job.

The rest of this case uses illustrative operating assumptions. \textbf{No customer, enrollment, fee, conversion rate, contract, measured learning gain, or margin described here should be construed as an accomplished Bayes result.}

\subsection*{77. First Choose the Unit of Service, Not a Duration}

Imagine an initial technical studio for working engineers who already possess a meaningful software foundation. 1 candidate target is the ability to design, evaluate, debug, and defend a production-minded AI feature on a bounded problem. The eventual program length should come from authentic task difficulty, readiness, practice cycles, available feedback, and learner sustainability: not from the desire to advertise a fixed number of months.

Before developing a flagship course, a smaller experiment might focus on 1 complete skill: diagnosing the reliability of a retrieval-based system using a baseline, error categories, and an improved evaluation. The experiment's unit is \textbf{a verified capability gain}, not ``4 weeks of content.''

This distinction affects economics immediately. A short pilot limits exposure and produces information about support minutes, student time, assessment difficulty, and the intensity of real demand. But it cannot automatically validate the price, motivation, or capacity requirements of a much longer course. The institute must earn that extrapolation through further evidence.

\subsection*{78. Decompose the First Offer into Six Production Functions}

\paragraph*{Acquisition and admission:} reach relevant working adults, explain the actual learning obligation, screen readiness, answer objections, and collect appropriate payment. A prospect who needs beginner instruction should not be placed into an advanced cohort merely to fill capacity.

\paragraph*{Learning design:} specify target capability, baseline tasks, prerequisite bridge, explanation, practice variants, and an independent end task. Some of this is reusable; some must change as the tools and audience change.

\paragraph*{Live instruction and facilitation:} teach, observe errors, answer consequential questions, and preserve a serious peer environment. These hours scale mostly with the number of groups, not exactly with learner count.

\paragraph*{Practice and feedback:} provide automated checks where valid, review individual attempts, diagnose substantive errors, and ensure revision. These hours are more variable and often the true capacity constraint.

\paragraph*{Assessment and evidence:} administer independent unfamiliar tasks, calibrate the rubric, handle borderline decisions, and report what is legitimately demonstrated. A marketing certificate is not a substitute.

\paragraph*{Operations and continuity:} schedule, provision, protect data, communicate policies, recover from incidents, manage refunds or accommodations, and capture lessons for the next iteration.

The cost model should correspond to these 6 functions, not just to a list of software subscriptions and instructor payments.

\subsection*{79. An Illustrative Weekly Founder-Hour Budget}

Suppose the founder can allocate 16 sustainable hours per week to Bayes for a short pilot. A possible planning budget is:

\begin{table*}[t]
\centering
\small
\caption{Illustrative Weekly Founder-Hour Budget for an Initial 16-Hour Delivery Capacity}
\begin{tabularx}{\textwidth}{p{0.34\textwidth}p{0.20\textwidth}X}
\toprule
\textbf{Work Category} & \textbf{Weekly Hours} & \textbf{Operational Notes} \\
\midrule
Live workshop / office hours & 3.0 & Combined interactive diagnosis and seminar delivery \\
Teaching preparation and updates & 2.0 & Falls only after reusable pedagogical assets stabilize \\
Admissions, sales, customer learning & 2.0 & Screening, readiness diagnostics, intake calls \\
Operations and coordination & 1.0 & Scheduling, platform integrity, quality monitoring \\
Senior assessment and standards & 1.0 & Uncompromised evaluation and rubric governance \\
Learner review and individual feedback & 5.0 & Highly capacity-sensitive; vulnerable to deadline peaks \\
Contingency and service recovery & 2.0 & Buffer protecting the customer promise and founder health \\
\midrule
\textbf{Total Sustainable Allocation} & \textbf{16.0} & \textbf{Binding constraint on early delivery capacity} \\
\bottomrule
\end{tabularx}
\end{table*}

This plan deliberately uses the entire available allotment and therefore has \textbf{no spare capacity outside the named contingency}. It is not a promise that 16 hours will suffice. Once measured work exceeds the plan, the institute must narrow scope, reduce enrollment, hire calibrated help, extend response times transparently, or stop accepting obligations. The founder's main job is not to force the workload into the estimate; it is to update the estimate from observed reality.

\subsection*{80. The Reviewer-Hour Problem Beneath an Apparently Simple Course}

Imagine 14 learners each submit 1 substantive design artifact every 2 weeks. Assume a skilled review averages 35 minutes, including reading, diagnosis, feedback, and recording the decision. Over an 8-week pilot, that is 4 submissions per learner: 56 reviews, totaling about \textbf{32.7 skilled hours} before rechecks, oral defenses, escalation, and grading calibration.

If every review requires the founder's judgment, average weekly review work is about \textbf{4.1 hours}: near the 5-hour budget above. But averages conceal deadline clustering. If most artifacts arrive on the same day, a 2-day turnaround may be infeasible even when the total weekly hours look adequate.

Possible design improvements include staged checkpoints, a clear project template that removes irrelevant setup, automated mechanical tests, a common debrief for errors affecting most learners, and reserved review windows. None justifies eliminating individual critique where that critique is the product's core value.

\subsection*{81. What Should Be Built Before the First Paid Cohort?}

Build only enough to ensure truthful delivery and useful learning. That might include a public sample lesson, a founder-authored example of strong and weak technical work, an entry diagnostic, a working project repository, a rubric, an individual submission route, and a simple scheduling and payment procedure. Use existing video, document, repository, and communication tools where they meet privacy and reliability requirements.

There is little reason to build a sophisticated custom LMS, record an entire 6-month curriculum, hire a full-time sales team, or develop an AI grading platform before discovering whether the target learners will pay and complete the 1st bounded experience.

The danger is not only financial. Heavy early infrastructure makes it psychologically and operationally harder to respond when the 1st buyers reveal a different urgent need.

\subsection*{82. The Asset Worth Preserving After a First Cohort}

After any pilot, preserve the \textbf{real production knowledge}: which prerequisites predicted difficulty, which tasks exposed genuine capability, where learner confusion clustered, how long substantive feedback took, which explanations worked, where AI assistance concealed gaps, and which promises caused misunderstandings.

Turn repeatable findings into a maintained error library, clearer exercises, revised rubric, faculty notes, and an explicit service-capacity model. Publish only suitable evidence with consent. The immediate economic benefit may be modest because the next cohort is small; the institutional benefit can be substantial if those assets let other teachers reproduce quality later.

A poorly documented 1st cohort is a service delivered once. A well-studied 1st cohort can be the beginning of an institutional production capability.
